\documentclass[12pt, a4paper]{article}

% Paquetes básicos
\usepackage[spanish]{babel}
\usepackage{amsmath, amssymb, amsfonts} % Para matemáticas avanzadas
\usepackage{geometry} % Para ajustar márgenes
\geometry{top=2.5cm, bottom=2.5cm, left=2.5cm, right=2.5cm}

% Datos del documento
\title{\textbf{Tarea 1 - Métodos Numéricos Avanzados}}
\author{Samuel Domínguez Gamero}
\date{Noviembre de 2026}

\begin{document}
	
\maketitle % Genera el título con los datos anteriores
\tableofcontents
\section{Tema 1}
\subsection{Adimensionalización de la ecuación}

\subsection{Existencia y Unicidad de la solución}
De acuerdo con la ecuación adimensionalizada, podemos realizar un análisis de la existencia y unicidad de las soluciones. En particular, el siguiente teorema resume las condiciones necesarias y suficientes para el análisis:

\begin{theorem}
\end{theorem}
\subsection{Resolución numérica: Método de Diferencias Finitas}
\subsection{Resolución numérica: Método del Disparo}

	
\end{document}
